% !TEX program = xelatex
% !TEX options = -synctex=1 -interaction=nonstopmode -file-line-error
% example-thesis.tex
% Demonstrates rossthesis: frontmatter, epigraph, chapter with sections,
% a table via rosstables, and bibliography.

\begin{filecontents}{example-thesis.bib}
@book{lezak2012neuropsych,
  author    = {Lezak, M. D. and Howieson, D. B. and Bigler, E. D.
               and Tranel, D.},
  title     = {Neuropsychological Assessment},
  edition   = {5},
  publisher = {Oxford University Press},
  year      = {2012},
}
@article{randolph1998rbans,
  author  = {Randolph, C.},
  title   = {RBANS manual},
  journal = {Psychological Corporation},
  year    = {1998},
}
\end{filecontents}

\documentclass[palette=oxford, font=palatino, bibstyle=authoryear]{rossthesis}
\addbibresource{example-thesis.bib}

% ── Thesis metadata ──────────────────────────────────────────────
\thesistitle{Neuropsychological Markers of Vascular Cognitive
             Impairment: A Multi-Method Investigation}
\thesisauthor{Ross A.\ Dunne}
\thesisdegree{Doctor of Philosophy}
\thesisuniversity{University of Limerick}
\thesisdepartment{Department of Psychology}
\thesisyear{2026}
\thesissupervisor{Prof.\ Ann-Marie Moriarty}

\begin{document}

% ── Frontmatter ──────────────────────────────────────────────────
\frontmatter
\maketitle
\declaration

\thesisabstract{%
  Vascular cognitive impairment (VCI) represents a heterogeneous spectrum
  of cognitive deficits attributable to cerebrovascular disease. This thesis
  examined the sensitivity of the RBANS to early VCI in three studies.
  Study~1 established normative data for Irish adults aged 50--90.
  Study~2 compared RBANS profiles across VCI subtypes. Study~3 investigated
  whether sleep architecture mediates the relationship between white-matter
  lesion burden and domain-specific cognition. Results support the utility
  of the RBANS as a screening instrument for VCI when combined with
  demographically adjusted norms.}

\acknowledgements{%
  I am deeply grateful to my supervisor, Prof.\ Ann-Marie Moriarty, for
  her unwavering guidance. Thanks also to the participants who generously
  gave their time, and to my family for their patience and support.}

\dedication{For my parents.}

\begin{abbreviations}
  \abbrev{RBANS}{Repeatable Battery for the Assessment of
                  Neuropsychological Status}
  \abbrev{VCI}{Vascular Cognitive Impairment}
  \abbrev{WML}{White-Matter Lesion}
  \abbrev{SWS}{Slow-Wave Sleep}
  \abbrev{PSG}{Polysomnography}
\end{abbreviations}

\makefrontlists

% ── Mainmatter ───────────────────────────────────────────────────
\mainmatter

\chapter{General Introduction}

\epigraph{The brain is wider than the sky.}{Emily Dickinson}

\section{Background}

Vascular cognitive impairment encompasses the full range of cognitive
deficits caused or exacerbated by cerebrovascular pathology
\parencite{lezak2012neuropsych}. Screening tools that are sensitive to
early VCI remain limited.

\section{The RBANS}

The Repeatable Battery for the Assessment of Neuropsychological Status
\parencite{randolph1998rbans} yields five index scores summarised in
\cref{tab:rbans-indices}.

\begin{table}[htbp]
  \centering
  \caption{RBANS index scores and contributing subtests.}
  \label{tab:rbans-indices}
  \begin{rosstable}{L{4cm} L{6cm} C{2.5cm}}
    \tablehead{Index & Subtests & Items} \\
    \midrule
    \shaderow Immediate Memory & List Learning, Story Memory & 30 \\
    Visuospatial/Constructional & Figure Copy, Line Orientation & 20 \\
    \shaderow Language & Picture Naming, Semantic Fluency & 20 \\
    Attention & Digit Span, Coding & 27 \\
    \shaderow Delayed Memory & List Recall, Story Recall, Figure Recall,
                                List Recognition & 24 \\
  \end{rosstable}
  \tablenote{Scores are age-adjusted with $M = 100$, $SD = 15$.}
\end{table}

\section{Aims of the Thesis}

This thesis aims to: (a)~establish Irish normative data for the RBANS;
(b)~characterise cognitive profiles across VCI subtypes; and
(c)~investigate sleep-mediated pathways to domain-specific impairment.

% ── Bibliography ─────────────────────────────────────────────────
\thesisbibliography

\end{document}
